% Generated by scripts/assemble.py; corrections belong in source/corrections.json.
\examhead{173}{西南财经大学 2024 年数学分析试卷}{164}
% p165-b05; source PDF page 165
\blockanchor{p165-b05}
\question{1}求极限 $\lim_{x\to \frac{\pi}{4}}\frac{\sec^2x - 2\tan x}{1 + \cos 4x}.$\par
% p165-b06; source PDF page 165
\blockanchor{p165-b06}
\question{2}设函数 $f(x)$ 在区间[0,1]上二阶可导，且 $f(0) = 0, f(1) = 1$.\par
% p165-b07; source PDF page 165
\blockanchor{p165-b07}
(1) 证明: 存在 $x_0 \in (0,1)$, 使得 $f'(x_0) = 1$.\par
% p165-b08; source PDF page 165
\blockanchor{p165-b08}
(2) 证明: 存在 $\xi \in (0,1)$, 使得 $\xi f''(\xi) + (1 + \xi)f'(\xi) = 1 + \xi$.\par
% p165-b09; source PDF page 165
\par\Needspace{6\baselineskip}
\blockanchor{p165-b09}
\question{3}设函数\par
\nopagebreak[4]
% p165-b10; source PDF page 165
\blockanchor{p165-b10}
\displayformula{f (x, y) = \left\{ \begin{array}{l l} x - y + \frac {x y ^ {3}}{x ^ {2} + y ^ {4}}, & (x, y) \neq (0, 0); \\ 0, & (x, y) = (0, 0). \end{array} \right.}
% p165-b11; source PDF page 165
\blockanchor{p165-b11}
讨论 $f(x,y)$ 在原点 $(0,0)$ 处是否连续，偏导数是否存在，是否可微.\par
% p165-b12; source PDF page 165
\par\Needspace{6\baselineskip}
\blockanchor{p165-b12}
\question{4}设 $z = z(x, y)$ 具有二阶连续偏导数, 请计算 $\frac{\partial^2 z}{\partial x^2} + \frac{\partial^2 z}{\partial y^2}$ 在坐标变换\par
\nopagebreak[4]
% p165-b13; source PDF page 165
\blockanchor{p165-b13}
\displayformula{\left\{ \begin{array}{l} u = x ^ {2} - y ^ {2}; \\ v = 2 x y. \end{array} \right.}
% p165-b14; source PDF page 165
\blockanchor{p165-b14}
下的表达式.\par
% p165-b15; source PDF page 165
\blockanchor{p165-b15}
\question{5}求 $f(x, y) = 2x^{2} + y^{2} - 8x - 2y + 9$ 在区域 $D = \{(x, y) \mid 2x^{2} + y^{2} \leq 1\}$ 上的最大值和最小值.\par
% p165-b16; source PDF page 165
\blockanchor{p165-b16}
\question{6}计算二重积分 $I = \iint_{D} |x - y| \, \mathrm{d}x \, \mathrm{d}y$，其中 $D = \{(x, y) \mid x^2 + y^2 \leq ax\}$，$a > 0$.\par
% p165-b17; source PDF page 165
\blockanchor{p165-b17}
\question{7}求下列级数的和 $\sum_{n=2}^{\infty} \frac{1}{2^n (n^2 - 1)}$.\par
% p165-b18; source PDF page 165
\blockanchor{p165-b18}
\question{8}判断反常积分 $\int_0^{+\infty} \frac{x^{p-1}}{1 + x} \, \mathrm{d}x$ 的敛散性.\par
